\chapter{Introduction}
\label{ch:intro}
Open with the motivation and the question your thesis answers. Set the
scene in a few paragraphs, then state your contributions explicitly.
Cross-references work throughout, for example Chapter~\ref{ch:method} and
Figure~\ref{fig:signal}.

\section{Background}
Survey the relevant prior work and position your contribution within it.
Cite sources with \texttt{\textbackslash cite}, for example a foundational
result~\cite{ref:turing} and a more recent study~\cite{ref:attention}.

\section{Contributions}
\begin{itemize}
  \item A clean separation of content and styling that scales to a full
        dissertation.
  \item A reproducible figure workflow rooted in the \texttt{figures/}
        directory.
  \item A minimalist cover that stays out of the reader's way.
\end{itemize}
